% This section is intended to be included via \input{} in a main paper file.
% It assumes common math packages such as amsmath, amssymb, and bm are already loaded.

\section{Methodology}
\label{sec:method}

\subsection{Problem Formulation}
\label{subsec:problem}

给定一段历史轨迹序列
\begin{equation}
\mathcal{X} = \{\bm{x}_t\}_{t=1}^{L}, \qquad \bm{x}_t \in \mathbb{R}^{C},
\end{equation}
其中 $\bm{x}_t$ 表示第 $t$ 个观测时刻的多维运动状态，包含位置、运动属性以及时间间隔等特征。我们的目标是在给定未来时间间隔序列
\begin{equation}
\bm{\delta} = \{\delta_h\}_{h=1}^{H}, \qquad \delta_h > 0,
\end{equation}
的条件下，预测未来 $H$ 个时刻的归一化二维运动增量
\begin{equation}
\hat{\mathcal{Y}}^{\mathrm{norm}} = \{\hat{\bm{y}}_h^{\mathrm{norm}}\}_{h=1}^{H}, \qquad \hat{\bm{y}}_h^{\mathrm{norm}} \in \mathbb{R}^{2},
\end{equation}
并进一步重建未来轨迹位置序列
\begin{equation}
\hat{\mathcal{P}} = \{\hat{\bm{p}}_h\}_{h=1}^{H}, \qquad \hat{\bm{p}}_h \in \mathbb{R}^{2}.
\end{equation}

与直接采用逐点递归预测不同，我们将整个模型构造为一个分层的 ``patch encoder + time-conditioned query decoder'' 框架。其核心思想是：首先在局部时间窗口内编码细粒度点级依赖，再在 patch 级别聚合长程历史上下文，最后依据未来时间间隔序列构造一组条件查询向量，通过交叉注意力从历史上下文中抽取与每个未来时刻最相关的信息。

\subsection{Patch-wise Hierarchical History Encoder}
\label{subsec:encoder}

\paragraph{Patch partition.}
我们将长度为 $L$ 的历史序列划分为 $N$ 个 patch，每个 patch 含有 $P$ 个观测点，即
\begin{equation}
L = N \times P.
\end{equation}
若输入长度与 $N \times P$ 不一致，则按照实现细节进行截断或零填充。记第 $n$ 个 patch 中第 $p$ 个点为 $\bm{x}_{n,p}$，则有
\begin{equation}
\mathcal{X}_n = \{\bm{x}_{n,p}\}_{p=1}^{P}, \qquad n = 1,\dots,N.
\end{equation}

为了显式编码不规则采样时间，我们根据历史中的时间间隔通道累积得到点级时间戳
\begin{equation}
\tau_{n,p} = \sum_{u=1}^{(n-1)P+p} \max(\Delta t_u, 0),
\end{equation}
并通过可学习时间映射 $\phi_t(\cdot)$ 生成时间嵌入
\begin{equation}
\bm{e}_{n,p} = \phi_t(\tau_{n,p}) \in \mathbb{R}^{d_t}.
\end{equation}
最终的点级输入特征写为
\begin{equation}
\bm{z}_{n,p} = [\bm{x}_{n,p}; \bm{e}_{n,p}] \in \mathbb{R}^{C + d_t}.
\end{equation}

\paragraph{Intra-patch self-attention.}
每个 patch 内部首先利用点级自注意力建模局部动态依赖。记第 $n$ 个 patch 的点特征矩阵为
\begin{equation}
\bm{Z}_n = [\bm{z}_{n,1}, \dots, \bm{z}_{n,P}]^\top \in \mathbb{R}^{P \times (C+d_t)}.
\end{equation}
我们先通过线性映射将其投影到隐藏空间，再施加多头自注意力与前馈网络：
\begin{align}
\bm{H}_n^{(0)} &= \bm{Z}_n \bm{W}_{\mathrm{in}}, \\
\tilde{\bm{H}}_n &= \mathrm{LN}\!\left(\bm{H}_n^{(0)} + \mathrm{MHA}\!\left(\bm{H}_n^{(0)}, \bm{H}_n^{(0)}, \bm{H}_n^{(0)}\right)\right), \\
\bm{H}_n^{(1)} &= \mathrm{LN}\!\left(\tilde{\bm{H}}_n + \mathrm{FFN}\!\left(\tilde{\bm{H}}_n\right)\right).
\end{align}
考虑到真实轨迹可能存在缺失观测，我们使用点级有效掩码 $m_{n,p} \in \{0,1\}$ 对 patch 内特征做掩码平均池化，从而得到局部 patch token：
\begin{equation}
\bm{h}_n^{\mathrm{loc}} =
\frac{\sum_{p=1}^{P} m_{n,p} \bm{H}_{n,p}^{(1)}}
{\max\!\left(\sum_{p=1}^{P} m_{n,p}, 1\right)}
\in \mathbb{R}^{d}.
\end{equation}

\paragraph{Patch existence bias.}
若第 $n$ 个 patch 中至少包含一个有效点，则记 patch 有效标记为 $a_n = 1$，否则 $a_n = 0$。我们引入一个可学习的 existence bias $\bm{b}_{\mathrm{ex}} \in \mathbb{R}^{d}$，用以显式区分 ``空 patch'' 与 ``有效 patch''：
\begin{equation}
\bm{h}_n^{(0)} = \bm{h}_n^{\mathrm{loc}} + a_n \bm{b}_{\mathrm{ex}}.
\end{equation}

\paragraph{Inter-patch temporal self-attention.}
局部 patch token 仅编码短时窗口内的动态模式，尚不足以建模长程依赖。为此，我们进一步在 patch 序列上施加多层上下文编码。与标准 Transformer 不同，我们在 patch 级注意力中加入显式的时间距离偏置，以鼓励模型优先关注时间上更接近的局部上下文。

对于第 $\ell$ 层的 patch token 序列 $\{\bm{h}^{(\ell-1)}_n\}_{n=1}^{N}$，任意两个 patch 之间的时间偏置定义为
\begin{equation}
b_{ij} = \log \left( \exp\left(- \frac{|i-j|\Delta_p}{\tau} \right) + \epsilon \right),
\end{equation}
其中 $\Delta_p$ 表示 patch 间的固定时间步长，$\tau$ 为衰减温度，$\epsilon$ 为数值稳定项。于是，带时间偏置的 patch 级注意力写为
\begin{equation}
\mathrm{Attn}_{ij}^{(\ell)} =
\frac{
\exp\!\left(
\frac{(\bm{q}_i^{(\ell)})^\top \bm{k}_j^{(\ell)}}{\sqrt{d_h}} + b_{ij}
\right)
}
{
\sum_{j'} \exp\!\left(
\frac{(\bm{q}_i^{(\ell)})^\top \bm{k}_{j'}^{(\ell)}}{\sqrt{d_h}} + b_{ij'}
\right)
},
\end{equation}
其中
\begin{equation}
\bm{q}_i^{(\ell)} = \bm{W}_Q \bm{h}_i^{(\ell-1)}, \quad
\bm{k}_j^{(\ell)} = \bm{W}_K \bm{h}_j^{(\ell-1)}, \quad
\bm{v}_j^{(\ell)} = \bm{W}_V \bm{h}_j^{(\ell-1)}.
\end{equation}
相应的聚合结果为
\begin{equation}
\bar{\bm{h}}_i^{(\ell)} =
\bm{h}_i^{(\ell-1)} +
\sum_{j=1}^{N} \mathrm{Attn}_{ij}^{(\ell)} \bm{v}_j^{(\ell)}.
\end{equation}

在时间偏置注意力之后，我们进一步叠加标准 Transformer 编码层以增强 patch 间的全局交互能力。若记位置编码为 $\mathrm{PE}(\cdot)$，则第 $\ell$ 层的上下文化输出为
\begin{equation}
\bm{h}_i^{(\ell)} =
\bar{\bm{h}}_i^{(\ell)} +
\mathrm{Transformer}^{(\ell)}\!\left(\mathrm{PE}\!\left(\bar{\bm{h}}_i^{(\ell)}\right)\right).
\end{equation}
经过 $M$ 层堆叠后，我们得到最终的历史上下文表示
\begin{equation}
\mathcal{H} = \{\bm{h}_1^{(M)}, \dots, \bm{h}_N^{(M)}\}.
\end{equation}

\subsection{Time-Conditioned Future Query Decoder}
\label{subsec:decoder}

\paragraph{Global historical anchor.}
为了使未来查询不仅依赖于时间间隔，还能够显式感知当前运动状态，我们首先构造一个历史全局锚点。具体地，记最后一个有效观测为 $\bm{x}_{\mathrm{last}}$，最后一个有效 patch 的上下文 token 为 $\bm{h}_{\mathrm{last}}^{(M)}$，则全局历史锚点定义为
\begin{equation}
\bm{g} = \psi_g([\bm{x}_{\mathrm{last}}; \bm{h}_{\mathrm{last}}^{(M)}]) \in \mathbb{R}^{d},
\end{equation}
其中 $\psi_g(\cdot)$ 为一个两层感知机。该设计将 ``当前观测状态'' 与 ``压缩后的历史上下文'' 统一到同一语义空间中，为后续未来查询提供稳定的条件起点。

\paragraph{Future interval embedding.}
对于未来第 $h$ 个预测步，我们首先对时间间隔做对数压缩
\begin{equation}
\tilde{\delta}_h = \log(1 + \delta_h),
\end{equation}
再通过可学习时间编码器 $\phi_f(\cdot)$ 得到其时间嵌入
\begin{equation}
\bm{d}_h = \phi_f(\tilde{\delta}_h) \in \mathbb{R}^{d_t'}.
\end{equation}
对数压缩有助于减弱大时间间隔带来的数值尺度不稳定问题，同时保留不同采样跨度之间的相对差异。

\paragraph{Step index embedding.}
若仅使用时间间隔构造 query，则当不同未来步具有相同或相近的间隔时，模型难以区分它们在预测序列中的语义位置。为此，我们进一步为未来第 $h$ 步引入一个可学习的步号嵌入
\begin{equation}
\bm{s}_h = \mathrm{Emb}(h) \in \mathbb{R}^{d}.
\end{equation}
这使得模型能够同时感知 ``多久以后'' 与 ``这是第几步'' 两类互补信息。

\paragraph{Future query construction.}
最终，第 $h$ 个未来步的条件查询向量由全局历史锚点、未来时间间隔嵌入以及步号嵌入共同构造：
\begin{equation}
\bm{q}_h = \psi_q([\bm{g}; \bm{d}_h; \bm{s}_h]) \in \mathbb{R}^{d},
\end{equation}
其中 $\psi_q(\cdot)$ 为用于 query 投影的多层感知机。由此，模型得到一组未来条件查询
\begin{equation}
\mathcal{Q} = \{\bm{q}_1, \dots, \bm{q}_H\}.
\end{equation}
这一设计本质上将未来预测转化为一个条件检索过程：每一个未来步都携带其自身的时间语义，并据此从历史 patch 表示中检索最相关的上下文信息。

\paragraph{Cross-attention decoding.}
给定未来查询序列 $\mathcal{Q}$ 和历史上下文序列 $\mathcal{H}$，我们采用多头交叉注意力完成历史到未来的信息路由：
\begin{equation}
\bm{c}_h = \mathrm{MHA}\!\left(\bm{q}_h, \mathcal{H}, \mathcal{H}\right),
\end{equation}
其中 $\bm{q}_h$ 为 query，$\mathcal{H}$ 同时作为 key 和 value。随后使用残差归一化得到未来条件表示
\begin{equation}
\tilde{\bm{c}}_h = \mathrm{LN}(\bm{c}_h + \bm{q}_h).
\end{equation}
最后，将交叉注意力输出、原始 query、最后观测以及时间间隔嵌入拼接，通过预测头回归归一化二维运动增量：
\begin{equation}
\hat{\bm{y}}_h^{\mathrm{norm}} = \psi_d([\tilde{\bm{c}}_h; \bm{q}_h; \bm{x}_{\mathrm{last}}; \bm{d}_h]) \in \mathbb{R}^{2},
\end{equation}
其中 $\psi_d(\cdot)$ 表示步级解码器。与单一全局解码头相比，上述设计允许不同未来时刻以不同的注意力模式访问历史上下文，因此更适合不规则时间间隔下的细粒度轨迹预测。

\subsection{Trajectory Reconstruction}
\label{subsec:reconstruction}

模型输出的是归一化后的运动增量 $\hat{\bm{y}}_h^{\mathrm{norm}}$。设运动尺度向量为 $\bm{s}_m \in \mathbb{R}^{2}$，位置尺度向量为 $\bm{s}_p \in \mathbb{R}^{2}$，则物理空间下的运动量为
\begin{equation}
\hat{\bm{y}}_h = \hat{\bm{y}}_h^{\mathrm{norm}} \odot \bm{s}_m.
\end{equation}

在速度建模模式下，未来第 $h$ 步的位移由速度与对应时间间隔相乘得到：
\begin{equation}
\Delta \hat{\bm{p}}_h = \hat{\bm{y}}_h \cdot \delta_h.
\end{equation}
在位移建模模式下，模型直接输出每一步的位移：
\begin{equation}
\Delta \hat{\bm{p}}_h = \hat{\bm{y}}_h.
\end{equation}
因此，未来轨迹位置通过累积和重建：
\begin{equation}
\hat{\bm{p}}_h = \sum_{r=1}^{h} \Delta \hat{\bm{p}}_r.
\end{equation}
对应的归一化位置为
\begin{equation}
\hat{\bm{p}}_h^{\mathrm{norm}} = \hat{\bm{p}}_h \oslash \bm{s}_p,
\end{equation}
其中 $\odot$ 和 $\oslash$ 分别表示按元素乘法与除法。

\subsection{Optimization Objective}
\label{subsec:loss}

我们从运动层面、完整轨迹层面以及最终落点层面联合优化模型。记真实未来运动增量为 $\bm{y}_h^{\mathrm{norm}}$，真实未来位置为 $\bm{p}_h^{\mathrm{norm}}$，则运动监督项定义为
\begin{equation}
\mathcal{L}_{\mathrm{motion}} =
\frac{1}{H} \sum_{h=1}^{H}
\ell\!\left(\hat{\bm{y}}_h^{\mathrm{norm}}, \bm{y}_h^{\mathrm{norm}}\right),
\end{equation}
其中 $\ell(\cdot,\cdot)$ 可取 Smooth L1 或 MSE 损失。

轨迹级重建损失写为
\begin{equation}
\mathcal{L}_{\mathrm{traj}} =
\frac{1}{H} \sum_{h=1}^{H}
\ell\!\left(\hat{\bm{p}}_h^{\mathrm{norm}}, \bm{p}_h^{\mathrm{norm}}\right),
\end{equation}
终点损失为
\begin{equation}
\mathcal{L}_{\mathrm{final}} =
\ell\!\left(\hat{\bm{p}}_H^{\mathrm{norm}}, \bm{p}_H^{\mathrm{norm}}\right).
\end{equation}
最终训练目标为三者的加权和：
\begin{equation}
\mathcal{L} =
\lambda_m \mathcal{L}_{\mathrm{motion}} +
\lambda_t \mathcal{L}_{\mathrm{traj}} +
\lambda_f \mathcal{L}_{\mathrm{final}},
\end{equation}
其中 $\lambda_m$、$\lambda_t$ 和 $\lambda_f$ 为权重系数。

\subsection{基于时间条件查询的局部 Patch 补全}
\label{subsec:completion}

除未来轨迹预测外，我们进一步考虑历史序列中的 patch 缺失补全问题。设第 $m$ 个 patch 为缺失 patch，其内部所有点均不可见，但其时间位置仍然已知。我们的目标是利用其局部时间邻域中的可见 patch 表示恢复该缺失 patch 的潜在表示，并进一步重建其内部轨迹点。与使用全局历史上下文做补全不同，我们采用固定大小的局部 chunk 进行补全，从而将模型归纳偏置显式约束在短时连续运动模式上。

\paragraph{Three-patch local chunk.}
对于每一个缺失 patch $m$，我们仅从其附近选择一个包含 $K=3$ 个可见 patch 的局部 chunk 作为补全上下文。记所选 patch 索引集合为
\begin{equation}
\mathcal{N}(m) = \{j_1, j_2, j_3\},
\end{equation}
其中每个 $j_k$ 对应于时间上邻近缺失 patch $m$ 的可见 patch，且满足 $a_{j_k}=1$。于是，局部补全记忆库写为
\begin{equation}
\mathcal{H}_m^{\mathrm{loc}} = \{\bm{h}_{j_1}^{(M)}, \bm{h}_{j_2}^{(M)}, \bm{h}_{j_3}^{(M)}\}.
\end{equation}
这种设计基于一个简单而有效的先验：对于短时轨迹序列，缺失片段的动态形态往往主要由其邻近上下文决定，而远距离历史信息更容易引入无关噪声。固定三 patch 局部窗口也使补全过程具备稳定的计算开销与明确的局部约束。

\paragraph{Missing-patch query construction.}
我们不直接将缺失 patch 视为一个待回归的匿名位置，而是为其构造一个具有时间语义、位置语义和局部状态语义的补全查询向量。设缺失 patch 的中心时间戳为 $\tau_m^{c}$，其 patch 索引嵌入为 $\bm{u}_m = \mathrm{Emb}(m)$。此外，我们从局部 chunk 中提取左右两侧最邻近可见 patch 的表示，分别记为 $\bm{h}_m^{-}$ 与 $\bm{h}_m^{+}$，并对整个局部 chunk 进行平均池化得到局部摘要
\begin{equation}
\bar{\bm{h}}_m = \frac{1}{3} \sum_{j \in \mathcal{N}(m)} \bm{h}_{j}^{(M)}.
\end{equation}
随后，缺失 patch 的条件 query 定义为
\begin{equation}
\bm{q}_m^{\mathrm{comp}} =
\psi_c\!\left(
\left[
\phi_c(\tau_m^{c});
\bm{u}_m;
\bm{h}_m^{-};
\bm{h}_m^{+};
\bar{\bm{h}}_m
\right]
\right),
\end{equation}
其中 $\phi_c(\cdot)$ 表示用于补全的时间编码器，$\psi_c(\cdot)$ 为 query 投影网络。该设计的关键在于：时间编码回答 ``缺失 patch 位于何时''，索引嵌入回答 ``缺失 patch 位于何处''，而左右邻域与局部摘要则回答 ``该处局部运动状态如何''。三者结合后，query 不再仅由时间戳单独决定，而能够显式表达局部补全所需的时空条件。

\paragraph{Local cross-attention completion.}
给定缺失 patch 查询 $\bm{q}_m^{\mathrm{comp}}$ 以及局部记忆库 $\mathcal{H}_m^{\mathrm{loc}}$，我们采用局部交叉注意力恢复其潜在 patch 表示：
\begin{equation}
\bm{r}_m = \mathrm{MHA}\!\left(\bm{q}_m^{\mathrm{comp}}, \mathcal{H}_m^{\mathrm{loc}}, \mathcal{H}_m^{\mathrm{loc}}\right),
\end{equation}
并通过残差归一化得到补全后的 patch token
\begin{equation}
\hat{\bm{h}}_m^{\mathrm{comp}} = \mathrm{LN}\!\left(\bm{r}_m + \bm{q}_m^{\mathrm{comp}}\right).
\end{equation}
由于 key 和 value 仅来自局部三 patch chunk，该补全过程可以被理解为一个局部条件检索器：模型依据缺失 patch 的时间与位置条件，从其紧邻的可见 patch 中自适应抽取最相关的动态模式，并将其映射到缺失 patch 的潜在语义空间。

\paragraph{Patch-level completion target.}
在训练阶段，我们将由完整输入编码得到的真实 patch 表示作为监督目标。记对应缺失 patch 的教师表示为 $\bm{h}_m^{\star}$，则 patch 级补全损失可写为
\begin{equation}
\mathcal{L}_{\mathrm{comp\_patch}} =
\ell_c\!\left(\hat{\bm{h}}_m^{\mathrm{comp}}, \bm{h}_m^{\star}\right),
\end{equation}
其中 $\ell_c(\cdot,\cdot)$ 可选用 MSE、Smooth L1 或 cosine 损失。该目标直接约束补全后的 patch token 与完整上下文下的目标表示保持一致，从而使补全过程首先在语义层面恢复缺失片段的整体动态结构。

\paragraph{Point-level patch reconstruction.}
在得到补全后的 patch token $\hat{\bm{h}}_m^{\mathrm{comp}}$ 后，我们进一步通过点级解码器恢复该 patch 内的细粒度轨迹点。记缺失 patch 内第 $p$ 个位置的相对时间编码为 $\phi_p(\tau_{m,p})$，则其点级重建结果写为
\begin{equation}
\hat{\bm{x}}_{m,p} = \psi_p\!\left([\hat{\bm{h}}_m^{\mathrm{comp}}; \phi_p(\tau_{m,p})]\right), \qquad p=1,\dots,P,
\end{equation}
其中 $\psi_p(\cdot)$ 表示点级重建网络。于是，点级补全损失定义为
\begin{equation}
\mathcal{L}_{\mathrm{comp\_point}} =
\frac{1}{P} \sum_{p=1}^{P}
\ell_p\!\left(\hat{\bm{x}}_{m,p}, \bm{x}_{m,p}^{\star}\right),
\end{equation}
其中 $\bm{x}_{m,p}^{\star}$ 为对应的真实点特征。

\paragraph{Joint objective for completion.}
最终，补全模块的训练目标写为
\begin{equation}
\mathcal{L}_{\mathrm{completion}} =
\lambda_{\mathrm{cp}} \mathcal{L}_{\mathrm{comp\_patch}} +
\lambda_{\mathrm{pt}} \mathcal{L}_{\mathrm{comp\_point}}.
\end{equation}
若将补全任务与主预测任务联合训练，则整体目标可进一步表示为
\begin{equation}
\mathcal{L}_{\mathrm{all}} =
\mathcal{L} + \lambda_{\mathrm{comp}} \mathcal{L}_{\mathrm{completion}},
\end{equation}
其中 $\mathcal{L}$ 表示前述未来轨迹预测损失，$\lambda_{\mathrm{comp}}$ 为补全分支权重。

上述补全机制与前述未来预测解码器在形式上是一致的，均采用 ``条件 query + cross-attention'' 的统一建模范式；但二者关注的上下文范围不同。未来预测依赖全局历史表示以支持长程时间外推，而缺失 patch 补全则被显式限制在局部三 patch chunk 内，以强化局部平滑性、连续性和结构一致性。该统一而分工明确的设计使模型能够同时处理历史补全与未来预测两类任务，并在不规则采样场景下保持稳定的时序归纳偏置。
